\documentclass[10pt,a4paper]{amsart}
\usepackage[margin=19mm]{geometry}
\usepackage{amsmath,amssymb,mathtools}
\usepackage[scr=rsfs,cal=euler]{mathalfa}
\usepackage{xfrac}
\usepackage[unicode,hidelinks,bookmarks=false]{hyperref}
\newcommand{\ie}{i.e.\ }
\newcommand{\cf}{cf.\ }
\newcommand{\resp}{respectively\ }
\newcommand{\Subsection}{\subsection}
\providecommand{\nobreakdash}{\nobreak\hbox{-}}
\setlength{\emergencystretch}{2em}
\multlinegap0pt
\usepackage{tikz}
\usetikzlibrary{cd}
\usepackage{amssymb}
\usepackage{mathtools}
\newtheorem{theorem}{Theorem}
\newtheorem{lemma}{Lemma}
\DeclareMathOperator{\Spec}{Spec}
\newcommand{\on}{\operatorname}
\title{On Moy-Prasad quotients over Laurent series fields}
\author{David Yang}
\date{Source: arXiv:2603.12623v1 (CC BY 4.0)}
\begin{document}
\maketitle

\noindent\emph{Source and licence.} On Moy-Prasad quotients over Laurent series fields. Version arXiv:2603.12623v1 (CC BY 4.0), distributed by arXiv under the Creative Commons Attribution 4.0 International licence, http://creativecommons.org/licenses/by/4.0/, as recorded in the arXiv licence metadata for this exact version and shown on the abstract page https://arxiv.org/abs/2603.12623v1. Full licence notice: \texttt{licenses/arxiv-2603.12623-CC-BY-4.0.txt}.\par\smallskip

\noindent\textit{Editorial scope.} Counted unit: the article's lemma inverseconstruction (source0.tex lines 848--859). The article's complete original proof of it is delivered with it (source0.tex lines 861--891, one \texttt{proof} environment of 31 lines). No mathematics is written, repaired or re-derived: the statement and the proof are the article's own lines, in the article's own notation, and no retained line is substituted or re-flowed.  Retained uncounted as the source-local context the proof needs: lines 386--399; lines 627--636.  Nothing was omitted from any retained span, so the package carries no omission record.  No retained line contains a citation, so the wrapper carries no bibliography and no bibliography entry.  The retained lines contain the article's own diagram code, which the wrapper typesets with the standard tikz packages; no image file is delivered and the package declares no asset dependency.  Editorial formatting only, in addition: an AMS \texttt{amsart} preamble that restates the article's own declarations and macros used by the retained lines, namely the environments \texttt{theorem}, \texttt{lemma} on separate counters, together with the standard packages those lines need.  The retained environments print in this document as Theorem 1, Lemma 1 in order of appearance, whereas the article prints the same items with the numbering of its own full document.  The complete native source of the article as distributed in the arXiv e-print for this exact version is bundled unchanged as \texttt{sources/196286/source0.tex}.
\begin{theorem}\label{depthcenter}
Choose $x\in\mathcal{B}(G)$ and a real number $r$.
\begin{enumerate}
\item The map $q$ sends $\mathfrak{k}_{x,r}$ into
\[
\displaystyle\bigoplus_{\substack{i\in\mathbb{Z}_+,j\in\mathbb{Q}\\ j\thinspace\geq\thinspace r\cdot i}}C_i((t))^j.
\]
\item The composition
\[
\mathfrak{k}_{x,r}\rightarrow\displaystyle\bigoplus_{\substack{i\in\mathbb{Z}_+,j\in\mathbb{Q}\\ j\thinspace\geq\thinspace r\cdot i}}C_i((t))^j\rightarrow\displaystyle\bigoplus_{i\in\mathbb{Z}_+}C_i((t))^{r\cdot i}
\]
factors through $\mathfrak{k}_{x,r}/\mathfrak{k}_{x,r+}.$
\end{enumerate}
\end{theorem}

\begin{lemma}\label{fqcompatibility}
There is a commutative diagram

\begin{center}
\begin{tikzcd}
\mathfrak{k}_{x,r}/\mathfrak{k}_{x,r+} \arrow[r, "q_{x,r}"]\arrow[d, "f_{x,r}"] &C((t))_{=r}\arrow[d]\\
\mathfrak{g}((t))\arrow[r, "q"] & C((t)).
\end{tikzcd}
\end{center}
\end{lemma}

\begin{lemma}\label{inverseconstruction}
Let $\overline{g}$ be any nonzero semisimple element of $\mathfrak{k}_{x,r}/\mathfrak{k}_{x,r+}$. Then:
\begin{enumerate}
\item There exists a semisimple element $g\in\mathfrak{k}_{x,r}$ lifting $\overline{g}$ such that $q(g)=q_{x,r}(\overline{g})$.
\item All choices of $g$ as in part (1) are $K_{x,0+}$-conjugate.
\item For any such $g$, we have $x\in\widetilde{B}(Z_G(g))$, and the equality
\begin{equation}\label{eq:centralizerdecomp}
\mathfrak{k}_{x,s}/\mathfrak{k}_{x,s+}\cong[\overline{g},\mathfrak{k}_{x,s-r}/\mathfrak{k}_{x,(s-r)+}]\oplus\mathfrak{k}^{Z_G(g)}_{x,s}/\mathfrak{k}^{Z_G(g)}_{x,s+}
\end{equation}
holds for any rational number $s$. Furthermore, $\mathfrak{k}^{Z_G(g)}_{x,s}/\mathfrak{k}^{Z_G(g)}_{x,s+}$ coincides with the kernel of $\on{ad}\overline{g}:\mathfrak{k}_{x,s}/\mathfrak{k}_{x,s+}\rightarrow\mathfrak{k}_{x,r+s}/\mathfrak{k}_{x,(r+s)+}$.
\end{enumerate}
\end{lemma}

\begin{proof}
First we show part (1). We can assume without loss of generality that $x\in\widetilde{A}(T).$ Then by Lemma \ref{fqcompatibility}, $g=f_{x,r}(\overline{g})$ satisfies the desired properties.

Next we show part (3), assuming part (2). Again assume that $x\in\widetilde{A}(T).$ By part (2), it suffices to show that the conditions of part (3) are satisfied for the specific choice $g=f_{x,r}(\overline{g})$. We note that it suffices to prove them after extending of $k((t))$, as the original statements can be recovered by taking the fixed points under the Galois action.

Because $(\mathfrak{k}_{x,r}/\mathfrak{k}_{x,r+})$ has a nonzero semisimple element, $r$ must be a rational number, and after potentially replacing $k((t))$ by an extension we can assume that $r$ is an integer and $G$ is split (so we can identify it with $H\times_{\Spec k}\Spec k((t))$). In this case, $\mathfrak{k}_{x,r}/\mathfrak{k}_{x,r+}$ contains $t^r\cdot\mathfrak{t}[[t]]/t^{r+1}\cdot\mathfrak{t}[[t]]$ as a Cartan subspace. Acting on $\overline{g}$ by an element of $L_x\subseteq G^x/K_{x,0+}$, we can assume that $\overline{g}\in t^r\cdot\mathfrak{t}[[t]]/t^{r+1}\cdot\mathfrak{t}[[t]].$ It then follows that $Z_G(f_{x,r}(\overline{g}))$ contains $T\times_{\Spec k}\Spec k((t))$, and hence, $\widetilde{\mathcal{B}}(Z_G(f_{x,r}(\overline{g})))$ contains $\widetilde{\mathcal{A}}(T)$, which contains $x$, as desired.

Let us show that (\ref{eq:centralizerdecomp}) holds in this case. Each term is a sum of affine root spaces. Let $\mathfrak{g}_{\alpha}t^m$ be an affine root space appearing in the decomposition of $\mathfrak{k}_{x,s}/\mathfrak{k}_{x,s+}$. Then it is easy to check that $\mathfrak{g}_{\alpha}t^m$ appears in $[\overline{g},\mathfrak{k}_{x,s-r}/\mathfrak{k}_{x,(s-r)+}]$ iff the element in $\mathfrak{t}$ corresponding to $\overline{g}$ does not pair with $\alpha$ to zero. On the other hand $\mathfrak{g}_{\alpha}t^m$ appears in $\mathfrak{k}^{Z_G(g)}_{x,s}/\mathfrak{k}^{Z_G(g)}_{x,s+}$ iff the value of this pairing is zero, so we get the desired direct sum decomposition. The same computation identifies $\mathfrak{k}^{Z_G(g)}_{x,s}/\mathfrak{k}^{Z_G(g)}_{x,s+}$ with the kernel of $\on{ad}\overline{g}.$

Finally, we treat part (2). Let $g$ be a lift of $\overline{g}$ satisfying the conditions in both part (1) and (3) (such an element was proven to exist above). Let $g_1$ be another lift of $\overline{g}$ satisfying the conditions of part (1). We claim that there is a $K_{x,0+}$-conjugate of $g_1$ which commutes with $g$. To see this, we construct an element $g(s)\in K_{x,0+}/K_{x,(s-r)+}$ for every real number $s\geq r$ with $\mathfrak{k}_{x,s}\neq\mathfrak{k}_{x,s+}$. This series of elements will satisfy two properties:
\begin{enumerate}
\item The $g(s)$-conjugate of the image of $g_1$ in $\mathfrak{k}_{x,r}/\mathfrak{k}_{x,s+}$ lies in $\mathfrak{k}^{Z_G(g)}_{x,r}/\mathfrak{k}^{Z_G(g)}_{x,s+}$.
\item For $s_2>s_1$, $g(s_1)$ is the image of $g(s_2)$ in $K_{x,0+}/K_{x,(s_1-r)+}$.
 \end{enumerate}
 Then when we conjugate $g_1$ by the inverse limit of the $g(s)$, we will get an element of $Z_{\mathfrak{g}}(g),$ as desired.

We can set $g(r)=1$. To construct the rest of the $g(s)$, we use induction. Let $s>r$ be a real number such that $\mathfrak{k}_{x,s}\neq\mathfrak{k}_{x,s+}$, and let $s_1$ be the largest number less than $s$ with $\mathfrak{k}_{x,s_1}\neq\mathfrak{k}_{x,s_1+}$. Assume that a value of $g(s_1)$ satisfying the properties has already been constructed. Then choose a lift $g(s)_0\in K_{x,0+}/K_{x,(s-r)+}$ lifting $g(s_1)$. By the inductive hypothesis, we will have
\[
\on{ad} g(s)_0(\overline{g}_1)\in\mathfrak{k}^{Z_G(g)}_{x,r}/\mathfrak{k}^{Z_G(g)}_{x,s+}+\mathfrak{k}_{x,s}/\mathfrak{k}_{x,s+},
\]
where $\overline{g}_1$ is the image of $g_1$ in $\mathfrak{k}_{x,r}/\mathfrak{k}_{x,s+}$. By part (3), the right hand side can be rewritten as $\mathfrak{k}^{Z_G(g)}_{x,r}/\mathfrak{k}^{Z_G(g)}_{x,s+}\oplus[\overline{g},\mathfrak{k}_{x,s-r}/\mathfrak{k}_{x,(s-r)+}]$. Let $g(s)_1$ be an element of $K_{x,s-r}/K_{x,(s-r)+}\cong\mathfrak{k}_{x,s-r}/\mathfrak{k}_{x,(s-r)+}$ such that $[\overline{g},g(s)_1]$ coincides with the image of $\on{ad} g(s)_0(\overline{g}_1)$ in $[\overline{g},\mathfrak{k}_{x,s-r}/\mathfrak{k}_{x,(s-r)+}]$. Then a short computation shows that $g(s)=g(s)_1g(s)_0$ satisfies the desired properties.

Let us replace $g_1$ with its $K_{x,0+}$-conjugate which commutes with $g$. We wish to show that $g_1=g.$ First, we claim that the images $c_1$ and $c$ of $g_1$ and $g,$ respectively, in $C^{Z_G(g)}\cong Z_{\mathfrak{g}}(g)//Z_G(g)$ agree. By Theorem \ref{depthcenter}, $c_1$ and $c$ are elements of $C^{Z_G(g)}_{\geq r}$ which are equal modulo $C^{Z_G(g)}_{>r}.$ Furthermore, by assumption, the images of $c_1$ and $c$ in $C$ are equal to $q_{x,r}(\overline{g})$.

It suffices to show that $c_1=c$ after an extension of $k((t))$, so we can assume that $r$ is an integer, $G$ and $Z_G(g)$ are split, and $Z_G(g)$ is a Levi subgroup of $G$. Then the bigradings on $C((t))$ and $C^{Z_G(g)}((t))$ with components $C_i((t))^j$ and $C^{Z_G(g)}_i((t))^j$, respectively, come from $\mathbb{G}_m\times\mathbb{G}_m$-actions, and the map $C^{Z_G(g)}((t))\rightarrow C((t))$ is $\mathbb{G}_m\times\mathbb{G}_m$-equivariant. Consider the inclusion $a:\mathbb{G}_m\rightarrow\mathbb{G}_m\times\mathbb{G}_m$ corresponding to the cocharacter $(r,-1)\in\mathbb{Z}\oplus\mathbb{Z}\cong X_*(\mathbb{G}_m\times\mathbb{G}_m).$ The fixed points of this action on $C((t))$ (resp. $C^{Z_G(g)}((t))$) are $C((t))_{=r}$ (resp. $C^{Z_G(g)}((t))_{=r}$.)

In particular, for any $z\in\mathbb{G}_m$, the images of $a(z)\cdot c_1$ and $c$ in $C((t))$ agree. However, if $c_1\neq c$, then $c_1$ cannot lie in $C^{Z_G(g)}((t))_{=r}$, and so this gives us infinitely many points of $C^{Z_G(g)}((t))$ which map to the same point of $C((t))$. But $C^{Z_G(g)}\rightarrow C$ is finite, so we get a contradiction, and $c_1$ must equal $c$.

Therefore, $g_1$ and $g$ must become $Z_G(g)$-conjugate after some extension of $k((t))$. But $g$ is central in $Z_G(g)$, so $g_1=g$, as desired.

\end{proof}

\end{document}
